\begin{table}[!htbp]\centering\footnotesize
\setlength{\tabcolsep}{4pt}
\caption{La transition année par année : l'État avance l'impôt, l'assiette le rembourse (Md€ courants, cas central)}\label{tab:is_transition}
\begin{tabular}{@{}l>{\raggedleft\arraybackslash}p{1.5cm}>{\raggedleft\arraybackslash}p{1.5cm}>{\raggedleft\arraybackslash}p{1.5cm}>{\raggedleft\arraybackslash}p{1.5cm}>{\raggedleft\arraybackslash}p{1.5cm}>{\raggedleft\arraybackslash}p{1.5cm}>{\raggedleft\arraybackslash}p{1.5cm}@{}}
\toprule
Année & Taux de déduction & Coût & dont permanent & Ressources & Solde & Dette en fin d'année & Écart de recette à 2026 \\
\midrule
2027 & 50~\% & 8,9 & 1,3 & 4,2 & $-$4,6 & 4,6 & $-$11,6 \\
2028 & 60~\% & 8,3 & 1,4 & 4,3 & $-$3,9 & 8,7 & $-$10,9 \\
2029 & 70~\% & 7,9 & 1,4 & 4,5 & $-$3,4 & 12,4 & $-$10,4 \\
2030 & 80~\% & 7,5 & 1,5 & 4,6 & $-$2,9 & 15,7 & $-$9,9 \\
2031 & 90~\% & 7,2 & 1,5 & 4,7 & $-$2,5 & 18,7 & $-$9,5 \\
2032 & 100~\% & 6,8 & 1,6 & 4,9 & $-$1,9 & 21,2 & $-$8,9 \\
2033 & 100~\% & 4,7 & 1,6 & 5,0 & 0,3 & 21,5 & $-$6,7 \\
2034 & 100~\% & 3,6 & 1,7 & 5,2 & 1,6 & 20,6 & $-$5,4 \\
2035 & 100~\% & 2,9 & 1,7 & 5,3 & 2,4 & 18,8 & $-$4,6 \\
2038 & 100~\% & 1,9 & 1,9 & 5,8 & 3,9 & 9,8 & $-$3,1 \\
2041 & 100~\% & 2,0 & 2,0 & 6,4 & 4,3 & éteinte & $-$2,7 \\
\multicolumn{8}{@{}l}{À partir de 2042 : baisse du taux normal de 1,0 point, qui rend l'excédent permanent aux entreprises.} \\
\bottomrule
\end{tabular}
\par\vspace{3pt}\parbox{\linewidth}{\raggedright\scriptsize Proposition : déduction immédiate de 100~\% des équipements neufs sous un plafond annuel de 1~M€ par groupe dès 2027 ; au-delà, 50~\% en 2027 puis 10 points de plus par an (taux de déduction au-delà du plafond). Coût de la déduction, ressources de l'assiette, solde de la réforme. Cas central : régime actuel pour moitié linéaire, pour moitié dégressif ; 75~\% de l'investissement éligible fait par des sociétés imposables et plafond couvrant 20~\% de cet investissement (hypothèses) ; ressources au milieu de la fourchette, en hausse de 3~\% par an ; intérêts à 3,25~\%. Coût permanent : ce que coûterait chaque année la déduction totale une fois la transition passée ; le reste du coût est une avance d'impôt, remboursée par les années suivantes. Écart de recette à 2026 : solde de la réforme moins les 7~Md€ de la contribution exceptionnelle perçus en 2026, qui s'éteint dans le droit en vigueur. La dette s'éteint en 2041 ; la baisse finale du taux (1,0 point) vient ensuite. Années d'exercice : en caisse, coûts et ressources arrivent environ un an plus tard. Source : \texttt{is/calculs\_note\_is.py}.}
\end{table}
